\subsection*{Висновки до першого розділу}
\addcontentsline{toc}{subsection}{Висновки до першого розділу}

У розділі окреслено теоретичні засади навчання майбутніх учителів інформатики розробки імерсивних освітніх ресурсів. Обґрунтовано, що створення та впровадження VR/AR/XR-рішень у навчальний процес виходить за межі суто технічної діяльності й потребує інтеграції технологічних, дидактичних і дизайнерських складників у структурі професійної підготовки.

Аналіз наукових джерел, структурований за п'ятьма проблемно-тематичними кластерами (загальнотеоретичний і дидактичний; освітні практики різних рівнів; професійна педагогічна підготовка; інклюзивний вимір; галузево-порівняльні підходи), засвідчив формування цілісного міждисциплінарного поля досліджень. Метааналітичні дані шести систематичних оглядів підтверджують педагогічну ефективність імерсивних технологій (розміри ефекту $g=0{,}82$--$d=0{,}98$), водночас когнітивно-афективна модель імерсивного навчання CAMIL \cite{Makransky2021CAMIL} визначає шість медіаторних факторів (присутність, агентність, втілення, когнітивне навантаження, мотивація, саморегуляція), що обумовлюють цю ефективність. Встановлено, що доповнена реальність зменшує зовнішнє когнітивне навантаження порівняно з традиційними засобами, тоді як віртуальна реальність може його збільшувати, що актуалізує роль дидактичного дизайну. Аналіз за моделлю TPACK виявив системний дефіцит TPK-компоненту (технолого-педагогічних знань): майбутні вчителі формують технологічні знання (TK) і педагогічні знання (PK) окремо, проте їх інтеграція в контексті імерсивних освітніх середовищ не забезпечується, що зумовлює підготовку вчителів переважно як \emph{користувачів}, а не \emph{розробників} імерсивних ресурсів.

Уточнення понятійного апарату дозволило визначити імерсивні освітні ресурси не як окремі цифрові продукти, а як інтерактивні дидактичні засоби нового покоління, побудовані на моделюванні віртуальних, доповнених і змішаних освітніх середовищ. Виокремлено п'ять сутнісних характеристик ІОР~-- інтерактивність, імерсивність/присутність, контекстуальність, візуалізація та адаптивність,~-- кожна з яких кореспондує з відповідними факторами моделі CAMIL. Порівняльний аналіз моделей педагогічного дизайну (ADDIE, SAM~I, LXD для AR) обґрунтував доцільність 6-компонентної моделі проєктування ІОР (аналіз~$\to$ вимоги~$\to$ вибір ПЗ~$\to$ розробка~$\to$ впровадження~$\to$ оцінювання), що вирізняється наявністю окремого етапу вибору програмного забезпечення та триетапною логікою навчання. Професійну готовність майбутнього вчителя інформатики до розробки ІОР визначено як інтегративну характеристику, що поєднує теоретичні знання, практичні вміння, особистісно-мотиваційні якості та здатність до безперервного професійного зростання. Наголошено, що педагогічна ефективність ІОР не реалізується автоматично, а залежить від якості дидактичного дизайну та готовності вчителя; ідентифіковано п'ять критеріїв ефективності (засвоєння знань, формування вмінь, мотивація, залученість, зручність), а також ризики~-- надмірна візуалізація без педагогічного обґрунтування, когнітивне перевантаження у VR, висока вартість обладнання, неготовність учителів.

Систематизація вітчизняного досвіду, здійснена в контексті міжнародних тенденцій, засвідчила, що підготовка вчителів переважно як користувачів, а не розробників імерсивних ресурсів є глобальним патерном, зафіксованим у систематичних оглядах. Аналіз нормативної рамки (проєкту Стандарту вищої освіти за спеціальністю 014) виявив, що компетентності ЗК17, ФК6, ФК11 та ПРН1--ПРН6 створюють лише технологічне підґрунтя, не артикулюючи педагогічного проєктування ІОР. Порівняльний аналіз ОПП 38~закладів вищої освіти (18~педагогічних, 14~класичних університетів, 6~спеціалізованих ЗВО) за запропонованою моделлю рівнів зрілості інтеграції ІОР показав, що абсолютна більшість перебуває на рівнях~0--2 (від повної відсутності імерсивних компонентів до інструментального використання в окремих дисциплінах), тоді як рівні~3--4 (проєктне створення ІОР та системна методика їх розробки) представлені одиничними випадками. Жодна з п'яти сутнісних характеристик ІОР не адресується як окремий об'єкт формування в проаналізованих ОПП. Дослідницька траєкторія КДПУ (2020--2024)~-- від порівняльного аналізу інструментів WebAR через проєктування квестів і формулювання 6-компонентної моделі до розроблення методики WebAR з інтеграцією машинного навчання~-- є єдиним задокументованим вітчизняним прикладом системного просування від рівня~2 до рівня~3--4, що засвідчує принципову можливість подолання виявленого розриву.

Виявлена суперечність між задекларованою значущістю імерсивних технологій в освітніх програмах і відсутністю системного підходу до формування відповідних професійних умінь емпірично обґрунтовує потребу в спеціально організованому педагогічному впливі. Методика навчання розробки імерсивних освітніх ресурсів постає не як факультативний компонент, а як необхідна умова формування готовності майбутнього вчителя інформатики до проєктування педагогічно доцільних і дидактично обґрунтованих цифрових середовищ.

Виявлені теоретичні положення та результати аналізу освітніх програм створюють підґрунтя для обґрунтування організаційно-педагогічних засад методики навчання розробки ІОР~-- дидактичних підходів і принципів, структурно-функціональної моделі та педагогічних умов, що розглядаються у наступному розділі.

